\clearpage\thispagestyle{empty}
\pdfbookmark[0]{Quellen zum Brief vom 25. August 1806}{zach18060825-sources}
\begin{center}\textbf{Quellen zum Brief vom 25. August 1806}\end{center}
\noindent Diese KI-Übertragung in heutiges Deutsch folgt dem Brief Carl Friedrich Gauß’ vom 25. August 1806 an Franz Xaver von Zach in \href{https://gdz.sub.uni-goettingen.de/id/PPN236007467}{Gauß’ \textit{Werken}, Band~VI} (Göttingen, 1874), S.~279--280. Der Brief erschien als Beitrag XXXIII in der \textit{Monatlichen Correspondenz}, Band~XIV, Oktober 1806, \href{https://www.digitale-sammlungen.de/en/view/bsb10538606?page=393}{S.~377}--\href{https://www.digitale-sammlungen.de/en/view/bsb10538606?page=398}{382}. Die Zeitschrift druckt die in der Werkausgabe beschriebene Ephemeride auf \href{https://www.digitale-sammlungen.de/en/view/bsb10538606?page=395}{S.~379}--\href{https://www.digitale-sammlungen.de/en/view/bsb10538606?page=397}{381} ab.
\par\medskip
Die Spalte für die Lichtstärke in der von Gauß angeführten Tabelle gibt für den 23. und 26. Februar 1806 den Wert 0,03983 und für den 1. März den Wert 0,03972 an. Im Brief steht für März der Wert 0.03983. \href{https://www.digitale-sammlungen.de/en/view/bsb10538603?page=492}{\textit{Monatliche Correspondenz}, Band~XI, Mai 1805, S.~478.}
\par\medskip
Im Anschluss an die Ephemeride druckt die Zeitschrift Gauß’ Berichtigungen zu seinem früheren Brief über den zweiten Kometen. Auf \href{https://www.digitale-sammlungen.de/en/view/bsb10538606?page=397}{S.~381} gibt er für S.~183, dritte Zeile von unten, \emph{Aberration} statt \emph{Nutation} an und für S.~186, vierte Zeile von unten, 71 statt 73. Auf \href{https://www.digitale-sammlungen.de/en/view/bsb10538606?page=398}{S.~382} erklärt er die Verkürzung um zwei Tage für bedeutungslos, da er in einem zweiten Beispiel für sein Werk mit ebenso gutem Erfolg Beobachtungen der Ceres im Abstand von 112 Tagen verwendet habe. Die historischen Digitalisate stammen von der Bayerischen Staatsbibliothek. Übertragung und Satz wurden mit KI erstellt.
\clearpage
